\documentclass[10.5pt,a4paper]{article}

% ============================================================
%  Currículo — Vinícius Cavati Gobbi
%  Compilar com: pdflatex cv-vinicius-gobbi.tex
%  (ou faça upload deste arquivo no Overleaf e clique em "Recompile")
% ============================================================

\usepackage[utf8]{inputenc}
\usepackage[T1]{fontenc}
\usepackage[brazilian]{babel}
\usepackage{lmodern}
\usepackage[margin=2cm,top=1.6cm,bottom=1.6cm]{geometry}
\usepackage[nobottomtitles]{titlesec}
\usepackage{enumitem}
\usepackage{xcolor}
\usepackage{fontawesome5}
\usepackage[hidelinks]{hyperref}
\usepackage{tabularx}
\usepackage{ragged2e}
\usepackage{fancyhdr}
\usepackage{needspace}

% ---------- Cores ----------
\definecolor{accent}{HTML}{0F766E}   % teal, alinhado à identidade visual do site
\definecolor{textgray}{HTML}{2B2B2B}
\definecolor{textmuted}{HTML}{5B5B5B}
\definecolor{ruleline}{HTML}{CCCCCC}

\color{textgray}
\setlength{\parindent}{0pt}
\setlength{\parskip}{0pt}

% ---------- Rodapé com número de página (só aparece se houver 2ª página) ----------
\pagestyle{fancy}
\fancyhf{}
\renewcommand{\headrulewidth}{0pt}
\fancyfoot[C]{\small\color{textmuted} \thepage}

% ---------- Seções ----------
\titleformat{\section}
  {\Large\bfseries\color{accent}}
  {}{0pt}{}
  [\vspace{2pt}{\color{ruleline}\hrule height 1pt}]
\titlespacing*{\section}{0pt}{16pt}{8pt}

% ---------- Listas com respiro ----------
\setlist[itemize]{leftmargin=16pt, topsep=3pt, itemsep=4pt, parsep=0pt}

% ---------- Comando para entradas (experiência / educação) ----------
% #1 título/empresa  #2 subtítulo/cargo  #3 período
\newcommand{\cventry}[3]{%
  \needspace{4\baselineskip}%
  \noindent
  \begin{tabularx}{\textwidth}{@{}X r@{}}
    \textbf{\large #1} & {\small\color{textmuted} #3} \\
  \end{tabularx}\\[1pt]
  {\itshape\color{accent} #2}
}

% ---------- Comando para projetos ----------
% #1 nome  #2 link (pode ser vazio)  #3 stack  #4 descrição
\newcommand{\cvproject}[4]{%
  \needspace{4\baselineskip}%
  \noindent
  \begin{tabularx}{\textwidth}{@{}X r@{}}
    \textbf{#1} & \ifx&#2&{}\else{\small\href{https://#2}{#2}}\fi \\
  \end{tabularx}\\[1pt]
  {\small\color{accent}\itshape #3}\\[2pt]
  {\small #4}
}

\begin{document}

% ============================================================
% CABEÇALHO
% ============================================================
\begin{center}
  {\Huge\bfseries\color{accent} Vinícius Cavati Gobbi}\\[5pt]
  {\large Estudante de Ciência da Computação \,\textbardbl\, Desenvolvedor Full-Stack}\\[8pt]
  {\small
    \faMapMarker\ Cariacica, ES -- Brasil
    \hspace{10pt}\textcolor{ruleline}{\vline}\hspace{10pt}
    \faEnvelope\ \href{mailto:vinicius.cgobbi2004@gmail.com}{vinicius.cgobbi2004@gmail.com}
    \hspace{10pt}\textcolor{ruleline}{\vline}\hspace{10pt}
    \faLinkedin\ \href{https://linkedin.com/in/vinicgobbi}{linkedin.com/in/vinicgobbi}\\[4pt]
    \faGithub\ \href{https://github.com/vinicgobbi}{github.com/vinicgobbi}
    \hspace{10pt}\textcolor{ruleline}{\vline}\hspace{10pt}
    \faGlobe\ \href{https://vinicgobbi.dev.br}{vinicgobbi.dev.br}
  }
\end{center}

% ============================================================
% RESUMO
% ============================================================
\section*{Resumo Profissional}
\justifying
Desenvolvedor Full-Stack apaixonado pelo ecossistema Linux e pela cultura Open-Source.
Atua atualmente na FAESA no desenvolvimento de sistemas com Laravel e React, com
experiência sólida também em Angular, o que proporciona versatilidade para transitar
entre diferentes stacks. Busca construir soluções que funcionem de verdade, com
interfaces intuitivas e código limpo e organizado.

% ============================================================
% EXPERIÊNCIA
% ============================================================
\section*{Experiência Profissional}

\cventry{FAESA}{Desenvolvedor Full-Stack (Estagiário)}{Jul 2025 -- Presente}
\begin{itemize}
  \item Atuação no desenvolvimento Full-Stack de sistemas acadêmicos e administrativos utilizando Laravel (PHP), React e JavaScript
  \item Modelagem de dados, gerenciamento de banco de dados SQL Server e criação de Stored Procedures para automação de rotinas e notificações
  \item Criação, manutenção e integração de APIs RESTful para comunicação eficiente entre sistemas internos
  \item Administração de chamados técnicos e suporte nível 2 via plataforma GLPI, garantindo a resolução de falhas em ambiente de produção
\end{itemize}

\vspace{10pt}
\cventry{Oi Alimentos}{Assistente / Auxiliar de PCM}{Nov 2022 -- Jul 2025}
\begin{itemize}
  \item Promovido de Jovem Aprendiz para Auxiliar e, posteriormente, para Assistente, demonstrando rápido aprendizado e adaptação
  \item Controle de Ordens de Serviço com o sistema Engeman
  \item Requisição de materiais com o sistema Agrosys
  \item Atendimento a Solicitações de Serviço baseado em prioridade
  \item Rotinas organizacionais através do método Kanban
\end{itemize}

% ============================================================
% EDUCAÇÃO
% ============================================================
\section*{Educação}
\cventry{FAESA}{Bacharelado em Ciência da Computação}{2023 -- 2026}

% ============================================================
% PROJETOS
% ============================================================
\section*{Projetos}

\cvproject{Amigo Indica}{amigoindica.faesa.br}
{Laravel · React · TypeScript · SQL Server}
{Sistema para gerenciar o programa de indicações da FAESA, facilitando a captação de novos alunos através de recomendações.}

\vspace{10pt}
\cvproject{Simulador do ENEM}{simuladorenem.faesa.br}
{Symfony}
{Simulador de descontos com base na nota do ENEM, com integração ao RD Station para captação de leads.}

\vspace{10pt}
\cvproject{Instrumentos Avaliativos}{}
{Laravel · React · TypeScript · SQL Server}
{Sistema para gerenciar notas e acompanhar o desempenho dos alunos de medicina da FAESA.}

\vspace{10pt}
\cvproject{Sistema de Agendamentos -- Psicologia}{}
{Laravel · SQL Server}
{Gestão de agendamentos e otimização do controle de horários para a clínica de Psicologia da FAESA.}

\vspace{10pt}
\cvproject{Sistema de Agendamentos -- Odontologia}{}
{Laravel · SQL Server}
{Gestão de agendamentos e otimização do controle de horários para a clínica de Odontologia da FAESA.}

\vspace{10pt}
\cvproject{Single Sign-On (SSO) -- AVA FAESA}{}
{Laravel · SAML}
{Simplificação do acesso dos alunos ao Ambiente Virtual de Aprendizagem, integrado à infraestrutura D2L via protocolo SAML.}

\vspace{10pt}
\cvproject{Notificações Automatizadas -- FAESA APP}{}
{SQL Server}
{Stored Procedures para automatizar o envio de notificações de faltas, atualizações de notas e pendências financeiras.}

\vspace{10pt}
\cvproject{Portfólio Pessoal}{vinicgobbi.dev.br}
{Angular · TypeScript}
{Cartão de visitas digital para expor habilidades, projetos e trajetória como desenvolvedor.}

% ============================================================
% STACKS / HABILIDADES TÉCNICAS
% ============================================================
\section*{Stacks e Ferramentas}
\begin{tabularx}{\textwidth}{@{}l @{\hspace{14pt}} X@{}}
  \textbf{Linguagens} & JavaScript, TypeScript, PHP, Python \\[7pt]
  \textbf{Frameworks \& Bibliotecas} & React, Angular, Laravel, Bootstrap \\[7pt]
  \textbf{Banco de Dados \& API} & MS SQL Server, MySQL, Postman \\[7pt]
  \textbf{DevOps \& Versionamento} & Docker, Git, GitHub \\[7pt]
  \textbf{Ambiente \& Ferramentas} & PNPM, VS Code, Linux, Windows, ZSH \\
\end{tabularx}

% ============================================================
% CERTIFICAÇÕES E CURSOS
% ============================================================
\section*{Certificações e Cursos}
\begin{itemize}
  \item \textbf{MikroTik Certified Network Associate (MTCNA)} -- MikroTik, 05/2026 \\
    {\small\color{textmuted}\itshape válido até 05/2029 \,\textbardbl\, código: 2605NA3984}
  \item \textbf{PRIVACIDADE E PROTEÇÃO DE DADOS (LGPD)} -- SENAI, 12/2022
  \item \textbf{Introdução à Computação (INTROCOMP)} -- UFES, 10/2021
  \item \textbf{Informática} -- Dat@School, 06/2017
  \item \textbf{Mecânico de Máquinas Industriais} -- SENAI, 06/2024
\end{itemize}

\end{document}
